\section{Methodology}
\label{sec:method}

We aim to measure the "addressability threshold" and its relationship to \reasoning in \gpt. Our approach involves presenting a model with a "Geometry Definition" followed by queries that test either its retrieval of that definition (\addressing) or its understanding of the spatial layout (\reasoning).

\para{Dataset.}
We use the \kilogram dataset~\cite{ji2022kilogram}, which consists of 1,018 tangram shapes. Each shape is composed of several parts, each defined by a set of 2D coordinates in an \svg-like format. For our experiments, we selected a subset of 20 distinct shapes to ensure variety in complexity and part arrangement.

\para{Experimental Design.}
We conduct three primary experiments:
\begin{enumerate}[leftmargin=*,itemsep=0pt,topsep=0pt]
    \item {\bf Experiment 1: \addressing Accuracy.} We provide a geometry definition consisting of part names and their corresponding coordinates. We then ask: "What are the coordinates of [Part Name]?" We vary the number of examples (shots) $N \in \{0, 1, 3, 5, 10\}$.
    \item {\bf Experiment 2: \reasoning Accuracy.} Using the same geometry definitions, we ask relative position queries such as "Is [Part A] above [Part B]?" The answers are binary (Yes/No). We vary $N \in \{1, 5, 10, 20\}$ shots.
    \item {\bf Experiment 3: Control Conditions.} We repeat the \reasoning tasks using (a) novel nonsense names for parts to eliminate semantic priors and (b) Chain-of-Thought (\cotask) prompting to test if explicit reasoning steps improve performance.
\end{enumerate}

\para{Model and Prompting.}
All experiments were conducted using \gpt. Prompts were structured to clearly separate the geometry definition from the task instructions. For \addressing, we used a robust extraction pipeline to handle minor formatting variations in the model's output, comparing the extracted coordinates to the ground truth. For \reasoning, we used exact match on the binary response.

\para{Evaluation Metrics.}
\begin{itemize}[leftmargin=*,itemsep=0pt,topsep=0pt]
    \item {\bf Addressing Accuracy:} The percentage of queries where the model returned the exact ground truth coordinate string for the requested part.
    \item {\bf Reasoning Accuracy:} The percentage of correct binary answers to relative position queries.
\end{itemize}
